\subsection{Compatibilité des moteurs et changements de règles}
\label{sec:compatibilite}
Trois niveaux de contrôle doivent être distingués : conserver les
fonctionnalités du laboratoire de simulation, reproduire une trajectoire
dans un régime de référence, et comparer deux modèles dont les règles
diffèrent. Seul le deuxième porte directement sur la parité dynamique.

Le rapport M4.2, section « Parité M4B », documente des tests à
$A=1$, $\gamma=1/2$ et $k=2$ : identité des variables discrètes
et comparaison des flottants à tolérance relative $10^{-9}$.
La sonde de 1000 pas rapportée ne présente pas de divergence discrète,
avec un écart relatif maximal des séries de $3{,}1\,10^{-12}$.
Ce constat n'est pas une identité bit à bit ni une garantie à tout horizon.

Pour M4.4, dans le dossier de campagne
\chemin{m4_4_rebond_credit_soc/}, le fichier
\chemin{results/analysis/parity_deviations_8000.csv}
conserve 8000 écarts maximaux nuls sur les 26 colonnes comparées à
la réalisation M4.3 de référence, dans le régime homogène à cible
arithmétique. Le test \chemin{tests/test_parity_m4_3.py} définit ce
périmètre. Le test \chemin{tests/test_v1_equivalence.py} compare
séparément Live-v1 au mode de compatibilité
\texttt{richest\_lends}, avec une intervention hétérogène ; un contrôle
négatif vérifie que libérer le sens du prêt peut changer la trajectoire.
La présence de ce test définit une exigence vérifiable ; elle ne remplace
pas à elle seule un compte rendu d'exécution pour chaque version.

Les chemins de campagnes M4.4 sont conservés comme identifiants de
provenance. Le passage à une règle de marché différente doit être
distingué d'une extension d'instrumentation, sans renommer les archives
dont proviennent les données. La non-régression ne
permet donc pas de fusionner les échantillons M4B, Live et M4.4.

\subsection{Pourquoi le regroupement temporel n'est pas une symétrie}
\label{sec:composition_temps}
Un contre-exemple apparaît avant même d'introduire le marché. Sans bruit,
crédit ni naissances, posons $F(K)=(1-\delta)(K+AK^\gamma)$.
Deux pas successifs donnent
\[
F(F(K))=(1-\delta)^2(K+AK^\gamma)
 +(1-\delta)A\big[(1-\delta)(K+AK^\gamma)\big]^\gamma.
\]
Un pas regroupé, avec $A'=2A$ et
$1-\delta'=(1-\delta)^2$, donnerait au contraire
\[
G(K)=(1-\delta)^2(K+2AK^\gamma).
\]
En général $F\circ F\ne G$ : la seconde production se calcule sur
un stock déjà modifié. Par exemple, à $K=100$, $A=1$,
$\gamma=1/2$ et $\delta=0{,}01$, les valeurs sont environ
$118{,}142$ et $117{,}612$. Le marché ajoute une alternance de
rencontres, de contrats, d'intérêts et de faillites ; composer les lois
de Poisson ou les seuls chocs ne compose pas ces opérations couplées.

Les essais Live-v1 de la figure~\ref{fig:art_temps} mesurent la portée
collective de cette différence. Après conversion complète des paramètres,
le regroupement de deux pas donne, relativement à sa référence,
$N:1{,}067\pm0{,}018$, $K_{\rm tot}:1{,}158\pm0{,}030$ et
$Y:1{,}118\pm0{,}024$. En regroupant quatre pas, le rapport de
production vaut $1{,}346\pm0{,}026$. Le protocole part d'un état vide,
utilise trois graines et compare $1200<t_{\rm ancien}\le2000$ ;
les flux sont divisés par le facteur de regroupement, pas les stocks.
Les valeurs par graine et les autres variantes sont dans
\chemin{data_article/temps_graines.csv} et
\chemin{data_article/temps_points.csv}. Ces essais et le contre-exemple
établissent l'absence d'une symétrie exacte ; ils ne constituent pas
une étude de convergence vers une limite continue.
